\documentclass[11pt]{amsart}
\usepackage[margin=1.15in]{geometry}
\usepackage{amsmath,amssymb,amsthm}
\usepackage[colorlinks=true,linkcolor=blue,citecolor=blue,urlcolor=blue]{hyperref}

\newtheorem{theorem}{Theorem}[section]
\newtheorem{lemma}[theorem]{Lemma}
\newtheorem{proposition}[theorem]{Proposition}
\newtheorem{corollary}[theorem]{Corollary}
\theoremstyle{definition}
\newtheorem{definition}[theorem]{Definition}
\theoremstyle{remark}
\newtheorem{remark}[theorem]{Remark}
\newtheorem{example}[theorem]{Example}

\newcommand{\N}{\mathbb{N}}
\newcommand{\Z}{\mathbb{Z}}
\newcommand{\Q}{\mathbb{Q}}
\newcommand{\R}{\mathbb{R}}
\newcommand{\cS}{\mathcal{S}}
\newcommand{\cT}{\mathcal{T}}
\newcommand{\cP}{\mathcal{P}}
\newcommand{\pb}{\bar p_5}
\newcommand{\Tr}{T}

\title[Two triangular numbers and a pentagonal number]{Every natural number is a sum of two triangular numbers and a pentagonal number}
\author{Quan Sun}
\date{\today}
\subjclass[2020]{Primary 11E25; Secondary 11D85, 11E20, 11P99}
\keywords{Triangular numbers, pentagonal numbers, universal sums, ternary quadratic forms, Eisenstein integers}

\begin{document}

\begin{abstract}
We prove that every nonnegative integer can be written as
$x(x+1)/2+y(y+1)/2+z(3z-1)/2$ with nonnegative integers $x,y,z$.
This confirms a conjecture of Zhi-Wei Sun from 2009.
More generally, let $P(x)$ be one of $x(3x-1)/2$, $x(3x+1)/2$, $x(3x+5)/2$ and let $d\ge 1$ be odd.
We show that every nonnegative integer has the form $P(x)+y(y+d)/2+z(z+1)/2$ with nonnegative integers $x,y,z$
whenever $d\le 17$, $d\le 7$, $d\le 3$ respectively; by a result of Z.-W. Sun these bounds cannot be improved.
We also show that every nonnegative integer is the sum of a triangular number,
a pentagonal number $y(3y-1)/2$ and a second pentagonal number $z(3z+1)/2$ with nonnegative $y,z$,
and that in this statement and in the one for $x(x+1)/2+y(y+1)/2+z(3z+1)/2$ one can take $z\ge 1$ as soon as the integer is at least $2$.
Finally, every positive integer is the sum of a pentagonal number and two triangular numbers with different indices.
Altogether this settles twenty-four of the $407$ conjectural universal tuples over $\N$ listed by Z.-W. Sun in 2020.
The proof starts from a primitive representation by $x^2+y^2+3z^2$ or $x^2+3y^2+3z^2$, which the classical
theory of genera provides. The new part is elementary: multiplication by a cube root of unity in the ring of
Eisenstein integers moves between representations, and a rotation of infinite order shows that along these moves
the sign condition behind the word ``nonnegative'' is eventually met.
For twenty-three of the twenty-four tuples this new part, and the deduction of the results from the classical input, have been formalized in Lean~4.
\end{abstract}

\maketitle

\section{Introduction}

Let $\N=\{0,1,2,\dots\}$. The triangular numbers are $\Tr(x)=x(x+1)/2$ $(x\in\N)$,
the pentagonal numbers are $p_5(x)=x(3x-1)/2$ $(x\in\N)$, and the numbers
$\pb(x)=p_5(-x)=x(3x+1)/2$ $(x\in\N)$ are called second pentagonal numbers.
More generally $p_m(x)=(m-2)x(x-1)/2+x$ denotes the $m$-gonal numbers.
Following Z.-W. Sun \cite{S20}, for integers $k\ge 1$ and $l>-k$ with $k\equiv l\pmod 2$ put
\[
\psi_{k,l}(x)=\frac{x(kx+l)}{2},
\]
so that $\Tr=\psi_{1,1}$, $p_5=\psi_{3,-1}$ and $\pb=\psi_{3,1}$.
A tuple $(k_1,l_1,k_2,l_2,k_3,l_3)$ is called \emph{universal over $\N$} if every $n\in\N$ can be written as
$\psi_{k_1,l_1}(x)+\psi_{k_2,l_2}(y)+\psi_{k_3,l_3}(z)$ with $x,y,z\in\N$, and \emph{universal over $\Z$}
if this holds with $x,y,z\in\Z$.

Universality over $\Z$ is a question about the integers represented by a ternary quadratic form
under congruence conditions, and a large part of the theory is available for it.
Universality over $\N$ adds sign conditions, which are not congruence conditions, and is much harder.
In \cite{S15} Z.-W. Sun showed that if $p_i+p_j+p_k$ $(3\le i\le j\le k,\ k\ge 5)$ is universal over $\N$,
then $(i,j,k)$ is one of $31$ explicit triples, and conjectured \cite[Conjecture~1.10]{S15}
that all $31$ sums are universal over $\N$; a second list of $64$ sums $ap_i+bp_j+cp_k$ is the subject of
\cite[Conjecture~1.13]{S15}. In \cite{S20} he showed that a universal tuple over $\N$ of a certain general shape must be one of $473$ candidates, proved $56$
of them universal over $\N$, and conjectured the same for the remaining $10+407$
\cite[Conjectures~1.2 and~1.5]{S20}. According to \cite{S20}, at that time none of the $95$ sums of \cite{S15}
had been proved universal over $\N$.
In 2026 Cao, Guo, Qiu, Feng and Gao \cite{CGQFG} proved the tuple $(5,1,3,1,1,1)$ universal over $\N$.
We are not aware of a proof of any of the statements below in the literature.

Our first result settles the triple $(3,3,5)$ of \cite[Conjecture~1.10]{S15}.

\begin{theorem}\label{thm:main}
Every $n\in\N$ can be written as
\[
n=\frac{x(x+1)}{2}+\frac{y(y+1)}{2}+\frac{z(3z-1)}{2}\qquad\text{with } x,y,z\in\N .
\]
If $n\ge 1$, one can moreover take $z\ge 1$.
\end{theorem}

Theorem~\ref{thm:main} is a special case of the following result on three families of tuples.

\begin{theorem}\label{thm:family}
Let $d$ be a positive odd integer.
\begin{enumerate}
\item[(i)] If $d\le 17$, the tuple $(3,-1,1,d,1,1)$ is universal over $\N$.
\item[(ii)] If $d\le 7$, the tuple $(3,1,1,d,1,1)$ is universal over $\N$.
\item[(iii)] If $d\le 3$, the tuple $(3,5,1,d,1,1)$ is universal over $\N$.
\end{enumerate}
In each case the tuple is not universal over $\N$ for larger odd $d$.
\end{theorem}

Thus, for instance, every $n\in\N$ is of the form $x(3x-1)/2+y(y+17)/2+z(z+1)/2$ with $x,y,z\in\N$,
while $9$ is not of the form $x(3x-1)/2+y(y+19)/2+z(z+1)/2$.
The last sentence of Theorem~\ref{thm:family} is not new: it is contained in \cite[Theorem~1.4(ii)]{S20},
and we include the one-line verification only for completeness.
The fifteen tuples of Theorem~\ref{thm:family} are exactly the tuples of these three shapes in the list
\cite[Appendix]{S20}.

Since $\{\Tr(y)+\Tr(z):y,z\in\N\}=\{y(y+1)+z^2:y,z\in\N\}$ (see Lemma~\ref{lem:TT}), the case $d=1$ gives:

\begin{corollary}\label{cor:square}
The tuples $(3,-1,2,2,2,0)$, $(3,1,2,2,2,0)$ and $(3,5,2,2,2,0)$ are universal over $\N$.
In particular $2p_3+p_4+p_5$ is universal over $\N$: every $n\in\N$ can be written as
$x(x+1)+y^2+z(3z-1)/2$ with $x,y,z\in\N$.
\end{corollary}

The sum $2p_3+p_4+p_5$ is one of the $64$ sums of \cite[Conjecture~1.13]{S15}.
Our third result concerns one triangular number and two pentagonal numbers of different kinds.

\begin{theorem}\label{thm:mixed}
Every $n\in\N$ can be written as
\[
n=\frac{x(x+1)}{2}+\frac{y(3y-1)}{2}+\frac{z(3z+1)}{2}\qquad\text{with } x,y,z\in\N .
\]
If $n\ge 1$, one can moreover take $y\ge 1$.
Consequently the two tuples $(3,1,3,-1,1,1)$ and $(3,5,3,1,1,1)$ are universal over $\N$.
\end{theorem}

The first assertion of Theorem~\ref{thm:mixed} was conjectured by Z.-W. Sun in the OEIS entry A254574 \cite{OEIS}.

In Theorem~\ref{thm:main} the pentagonal number can be taken positive, and in Theorem~\ref{thm:mixed} the same holds for $y(3y-1)/2$.
For the second pentagonal numbers $z(3z+1)/2$ the corresponding statements need a further argument.

\begin{theorem}\label{thm:shift}
Let $n\ge 2$ be an integer.
\begin{enumerate}
\item[(i)] There are $x,y\in\N$ and an integer $z\ge 1$ with $n=\dfrac{x(x+1)}{2}+\dfrac{y(y+1)}{2}+\dfrac{z(3z+1)}{2}$.
\item[(ii)] There are $x,y\in\N$ and an integer $z\ge 1$ with $n=\dfrac{x(x+1)}{2}+\dfrac{y(3y-1)}{2}+\dfrac{z(3z+1)}{2}$.
\end{enumerate}
Consequently the tuples $(3,7,1,1,1,1)$, $(3,7,2,2,2,0)$ and $(3,7,3,-1,1,1)$ are universal over $\N$.
\end{theorem}

The bound $n\ge 2$ cannot be improved, since $z(3z+1)/2\ge 2$ for $z\ge 1$.
Our last result is a refinement of Theorem~\ref{thm:main} in a different direction.

\begin{theorem}\label{thm:distinct}
Every positive integer $n$ can be written as
\[
n=\frac{x(3x-1)}{2}+\frac{r(r+1)}{2}+\frac{s(s+1)}{2}\qquad\text{with } x,r,s\in\N \text{ and } r\ne s .
\]
Consequently the tuple $(3,-1,2,4,2,2)$ is universal over $\N$: every $n\in\N$ can be written as $x(3x-1)/2+y(y+2)+z(z+1)$ with $x,y,z\in\N$.
\end{theorem}

Theorems~\ref{thm:family}, \ref{thm:mixed}, \ref{thm:shift} and~\ref{thm:distinct} and Corollary~\ref{cor:square} together settle
$15+3+2+3+1=24$ of the $407$ tuples of \cite[Conjecture~1.5]{S20}.

\subsection*{Method}
The identities $24\,\Tr(x)+3=3(2x+1)^2$, $24\,p_5(z)+1=(6z-1)^2$ and $24\,\pb(z)+1=(6z+1)^2$
turn Theorem~\ref{thm:main} into the statement that for every $n$ the equation
\begin{equation}\label{eq:intro}
24n+7=v^2+3a^2+3b^2
\end{equation}
has a solution in odd positive integers with $v\equiv 5\pmod 6$ or $v=1$.
Solutions in odd integers with $3\nmid v$ exist by the theory of ternary forms;
the difficulty is the residue of $|v|$ modulo $6$, that is, the sign of $v$ once $v$ is normalised by $v\equiv 1\pmod 3$.

The binary form $3a^2+v^2$ is four times the norm of $(v+a\sqrt{-3})/2$ in the ring $\Z[\omega]$
of Eisenstein integers, and multiplication by a cube root of unity produces from an odd pair $(a,v)$ exactly one
other odd pair $(a',v')$ with the same norm (Lemma~\ref{lem:partner}).
This move is not new: it is used for parity conditions in \cite[Lemma~3.2 and Theorem~1.7(iii)]{S15},
and it is one of the three moves of \cite[Lemmas~8 and~9]{CGQFG}.
The move either changes the residue of $v$ modulo $3$, or it is one explicit linear map $A$.
The same holds in the plane of $(b,v)$, with a linear map $B$.
What we add is the following observation.
If the residue of $v$ never changed along the alternating sequence of moves, the representations
$(BA)^j s_0$ $(j\ge 0)$ would all be integral and, being finite in number, periodic.
But $BA$ is a rotation of infinite order whose axis is the line $v=a=b$, and a $2$-adic argument
shows that only points of the axis have an integral two-sided orbit (Lemma~\ref{lem:2adic}).
Hence the residue changes, unless the starting point is $(k,k,k)$.
The remaining input is a starting point off the axis, which is supplied by the fact that
$x^2+3y^2+3z^2$ represents $24n+7$ primitively.
The side conditions $y(y+d)/2$ with $d>1$ are handled by a second, simpler observation:
when $v$ is large compared with $a$ and $b$, two consecutive moves both change the residue of $v$
and make $a$ large (Lemma~\ref{lem:large}); this leaves finitely many $n$, at most $779$, which are checked directly.
For Theorem~\ref{thm:shift} one has to avoid the value $v=1$ in addition.
Here the observation is that the periodicity argument still applies if a coordinate equal to $1$ is replaced by $-1$:
with this sign the moves that do not lead to a representation of the required kind are again given by the same two linear maps
(Lemma~\ref{lem:hat}).
In Theorem~\ref{thm:distinct} the representations with $a=b$ have to be avoided. Here the periodicity argument is not sufficient.
We show that a sequence of moves between two such representations has at most five steps,
unless it is excluded by a $2$-adic argument (Lemma~\ref{lem:val}); what remains is a short list of linear equations in two unknowns.

Excluding a periodic orbit of a rational isometry by showing that the relevant eigenvalues are not roots of unity
is also done in \cite[Lemma~6]{CGQFG}, there in order to pass from one class of a genus to the other.
In \cite{CGQFG} two independent sign conditions have to be met, and the argument ends with a computer-verified finite cover of a cone.
In the cases treated here only one binary condition is in play, and the periodicity argument settles it directly.

\subsection*{Organisation}
Section~\ref{sec:partner} introduces the partner map.
Section~\ref{sec:descent} contains the theorems on changing the residue and the lemma for large $v$.
Section~\ref{sec:seed} supplies primitive representations.
The theorems are proved in Section~\ref{sec:proofs}, except for Theorem~\ref{thm:distinct}, which is proved in Section~\ref{sec:distinct}.
Section~\ref{sec:lean} describes the formalization in Lean,
and Section~\ref{sec:remarks} contains remarks on the scope of the method and on the accompanying computations.

\section{The partner map}\label{sec:partner}

Let
\[
\cP=\{(u,v)\in\Z^2:\ u,v>0,\ u,v \text{ odd},\ 3\nmid v\}.
\]

\begin{definition}\label{def:partner}
For $(u,v)\in\cP$ define $\varphi(u,v)=(u',v')$ by
\[
(u',v')=
\begin{cases}
\bigl(\tfrac{u+v}{2},\ \tfrac{3u-v}{2}\bigr) & \text{if } u\equiv v\pmod 4 \text{ and } v<3u,\qquad \text{(case I)}\\[4pt]
\bigl(\tfrac{u+v}{2},\ \tfrac{v-3u}{2}\bigr) & \text{if } u\equiv v\pmod 4 \text{ and } v>3u,\qquad \text{(case II)}\\[4pt]
\bigl(\tfrac{|u-v|}{2},\ \tfrac{3u+v}{2}\bigr) & \text{if } u\not\equiv v\pmod 4.\qquad\qquad\qquad\ \ \text{(case III)}
\end{cases}
\]
The case $v=3u$ does not occur since $3\nmid v$.
\end{definition}

\begin{lemma}\label{lem:partner}
Let $(u,v)\in\cP$ and $(u',v')=\varphi(u,v)$. Then
\begin{enumerate}
\item[(a)] $(u',v')\in\cP$ and $3u'^2+v'^2=3u^2+v^2$;
\item[(b)] $v'\equiv v\pmod 3$ in case~I, and $v'\equiv -v\pmod 3$ in cases~II and~III;
\item[(c)] $\varphi(u',v')=(u,v)$.
\end{enumerate}
\end{lemma}

\begin{proof}
(a) For either choice of sign,
\[
3\Bigl(\frac{u\pm v}{2}\Bigr)^2+\Bigl(\frac{3u\mp v}{2}\Bigr)^2=\frac{12u^2+4v^2}{4}=3u^2+v^2 .
\]
Since $u$ and $v$ are odd, $u\equiv v\pmod 4$ or $u\equiv -v\pmod 4$.
If $u\equiv v\pmod 4$, then $u+v\equiv 2u\equiv 2\pmod 4$ and $3u-v\equiv 2u\equiv 2\pmod 4$,
so $(u+v)/2$ and $|3u-v|/2$ are odd; they are positive because $v\ne 3u$.
If $u\equiv -v\pmod 4$, then $u-v\equiv 2u\equiv 2\pmod 4$ and $3u+v\equiv 2u\equiv 2\pmod 4$,
so $|u-v|/2$ and $(3u+v)/2$ are odd and positive.
Finally $2v'\equiv \pm v\pmod 3$ shows $3\nmid v'$.

(b) As $2^{-1}\equiv -1\pmod 3$, we get $v'=(3u-v)/2\equiv v$ in case~I,
$v'=(v-3u)/2\equiv -v$ in case~II and $v'=(3u+v)/2\equiv -v$ in case~III.

(c) In case~I we have $u'+v'=2u\equiv 2\pmod 4$, so $u'\equiv v'\pmod 4$, and
$3u'-v'=2v>0$; hence $(u',v')$ is in case~I and $\varphi(u',v')=\bigl(\tfrac{u'+v'}{2},\tfrac{3u'-v'}{2}\bigr)=(u,v)$.
In case~II we have $u'-v'=2u\equiv 2\pmod 4$, so $u'\not\equiv v'\pmod 4$, and
$\varphi(u',v')=\bigl(\tfrac{|u'-v'|}{2},\tfrac{3u'+v'}{2}\bigr)=(u,v)$.
In case~III suppose first $u>v$. Then $u'+v'=2u\equiv 2\pmod 4$, so $u'\equiv v'\pmod 4$, and
$v'-3u'=2v>0$; hence $(u',v')$ is in case~II and
$\varphi(u',v')=\bigl(\tfrac{u'+v'}{2},\tfrac{v'-3u'}{2}\bigr)=(u,v)$.
If $u<v$, then $v'-u'=2u\equiv 2\pmod 4$, so $(u',v')$ is in case~III, and
$\varphi(u',v')=\bigl(\tfrac{v'-u'}{2},\tfrac{3u'+v'}{2}\bigr)=(u,v)$.
\end{proof}

Part~(c) justifies the name ``partner''; it is used only in Section~\ref{sec:distinct}.

\begin{remark}\label{rem:eisenstein}
Let $\omega=(-1+\sqrt{-3})/2$. For odd $u,v$ the number $\xi=(v+u\sqrt{-3})/2$ lies in $\Z[\omega]$ and has norm $(3u^2+v^2)/4$.
Write $\omega^j\xi=(v_j+u_j\sqrt{-3})/2$ for $j=0,1,2$. Then exactly one of the pairs $(u_j,v_j)$ consists of even numbers and
the other two consist of odd numbers, and $\varphi$ sends $(|u_j|,|v_j|)$ for one of the latter two to the same expression for the other.
Moreover $v_0+v_1+v_2=0$ and $v_0\equiv v_1\equiv v_2\pmod 3$.
The counting statement behind this picture is \cite[Lemma~3.2]{S15}.
\end{remark}

\section{Changing the residue}\label{sec:descent}

For a positive integer $N$ let
\begin{align*}
\cS(N)&=\{(v,a,b)\in\Z^3:\ v,a,b>0 \text{ odd},\ 3\nmid v,\ v^2+3a^2+3b^2=N\},\\
\cT(N)&=\{(u,v,w)\in\Z^3:\ u,v,w>0 \text{ odd},\ 3\nmid vw,\ 3u^2+v^2+w^2=N\}.
\end{align*}
Both sets are finite. We call a triple \emph{diagonal} if its three entries are equal.

\begin{lemma}\label{lem:2adic}
Let $M$ be one of the matrices
\[
M_1=\frac14\begin{pmatrix}5&-6\\ 6&-4\end{pmatrix},\qquad
M_2=\frac14\begin{pmatrix}-3&-2\\ 2&-4\end{pmatrix}.
\]
Then $\det M=1$, and if $y\in\Z^2$ satisfies $M^jy\in\Z^2$ for all $j\in\Z$, then $y=0$.
\end{lemma}

\begin{proof}
One checks $\det M_1=(-20+36)/16=1$ and $\det M_2=(12+4)/16=1$, so
\[
M_1^{-1}=\frac14\begin{pmatrix}-4&6\\ -6&5\end{pmatrix},\qquad
M_2^{-1}=\frac14\begin{pmatrix}-4&2\\ -2&-3\end{pmatrix}.
\]
Let $\Lambda=\{y\in\Z^2: M^jy\in\Z^2 \text{ for all } j\in\Z\}$. This is a subgroup of $\Z^2$ with $M\Lambda=\Lambda$.
We claim $\Lambda\subseteq 2\Z^2$. Let $y=(y_1,y_2)\in\Lambda$.

For $M=M_1$: the first entry of $M_1^{-1}y$ is $(-4y_1+6y_2)/4$, so $y_2$ is even; the first entry of $M_1y$ is $(5y_1-6y_2)/4$,
so $5y_1\equiv 6y_2\equiv 0\pmod 4$ and $y_1$ is even.

For $M=M_2$: the second entry of $M_2y$ is $(2y_1-4y_2)/4$, so $y_1$ is even; the first entry of $M_2^{-1}y$ is $(-4y_1+2y_2)/4$, so $y_2$ is even.

Now let $y\in\Lambda$. For every $j$ we have $M^jy\in\Lambda\subseteq 2\Z^2$, hence $M^j(y/2)\in\Z^2$, so $y/2\in\Lambda$.
Thus $\Lambda\subseteq 2\Lambda$, and by induction $\Lambda\subseteq 2^m\Z^2$ for all $m$. Therefore $\Lambda=\{0\}$.
\end{proof}

\begin{theorem}\label{thm:A}
Let $N$ be a positive integer. If $\cS(N)$ contains a triple that is not diagonal,
then $\cS(N)$ contains a triple $(v,a,b)$ with $v\equiv 1\pmod 3$ and a triple $(v,a,b)$ with $v\equiv 2\pmod 3$.
\end{theorem}

\begin{proof}
For $s=(v,a,b)\in\cS(N)$ put
\[
\alpha(s)=(v',a',b)\ \text{ with } (a',v')=\varphi(a,v),\qquad
\beta(s)=(v'',a,b'')\ \text{ with } (b'',v'')=\varphi(b,v).
\]
By Lemma~\ref{lem:partner}(a), $\alpha$ and $\beta$ map $\cS(N)$ to itself.
Let $s_0=(v_0,a_0,b_0)\in\cS(N)$ be not diagonal and define $s_i=(v_i,a_i,b_i)$ by
$s_{i+1}=\alpha(s_i)$ for even $i$ and $s_{i+1}=\beta(s_i)$ for odd $i$.
It suffices to show that $v_i\not\equiv v_0\pmod 3$ for some $i$.

Suppose, on the contrary, that $v_i\equiv v_0\pmod 3$ for all $i\ge 0$.
Since $3\nmid v_i$, Lemma~\ref{lem:partner}(b) shows that every step is in case~I, that is, $s_{i+1}=As_i$ for even $i$ and $s_{i+1}=Bs_i$ for odd $i$,
where $A$ and $B$ are the linear maps
\[
A(v,a,b)=\Bigl(\frac{3a-v}{2},\ \frac{a+v}{2},\ b\Bigr),\qquad
B(v,a,b)=\Bigl(\frac{3b-v}{2},\ a,\ \frac{b+v}{2}\Bigr).
\]
One checks $A^2=B^2=\mathrm{id}$, so $R=BA$ is invertible; explicitly
\[
R(v,a,b)=\Bigl(\frac{v-3a+6b}{4},\ \frac{v+a}{2},\ \frac{-v+3a+2b}{4}\Bigr).
\]
We have $s_{2j}=R^js_0$ for all $j\ge 0$. As $\cS(N)$ is finite, there are $0\le i<j$ with $R^is_0=R^js_0$,
hence $R^rs_0=s_0$ with $r=j-i\ge 1$.
Consequently $R^js_0\in\{s_0,Rs_0,\dots,R^{r-1}s_0\}\subseteq\Z^3$ for every $j\in\Z$.

Let $\pi(v,a,b)=(a-v,\ b-v)$. A direct computation from the formula for $R$ gives $\pi\circ R=M_1\circ\pi$:
indeed, writing $(y_1,y_2)=(a-v,b-v)$,
\[
\frac{v+a}{2}-\frac{v-3a+6b}{4}=\frac{5y_1-6y_2}{4},\qquad
\frac{-v+3a+2b}{4}-\frac{v-3a+6b}{4}=\frac{6y_1-4y_2}{4}.
\]
Hence $y=\pi(s_0)\in\Z^2$ satisfies $M_1^jy=\pi(R^js_0)\in\Z^2$ for all $j\in\Z$, and Lemma~\ref{lem:2adic} gives $y=0$,
that is, $a_0=v_0=b_0$. This contradicts the choice of $s_0$.
\end{proof}

\begin{theorem}\label{thm:B}
Let $N$ be a positive integer. If $\cT(N)$ contains a triple that is not diagonal,
then $\cT(N)$ contains a triple $(u,v,w)$ with $v\not\equiv w\pmod 3$.
\end{theorem}

\begin{proof}
We use the letters $\alpha,\beta,A,B,R,\pi$ for the objects analogous to those in the previous proof.
For $s=(u,v,w)\in\cT(N)$ put
\[
\alpha(s)=(u',v',w)\ \text{ with } (u',v')=\varphi(u,v),\qquad
\beta(s)=(u'',v,w'')\ \text{ with } (u'',w'')=\varphi(u,w);
\]
these are maps $\cT(N)\to\cT(N)$ by Lemma~\ref{lem:partner}(a).
Let $s_0\in\cT(N)$ be not diagonal and define $s_i=(u_i,v_i,w_i)$ by $s_{i+1}=\alpha(s_i)$ for even $i$ and $s_{i+1}=\beta(s_i)$ for odd $i$.
Suppose that $v_i\equiv w_i\pmod 3$ for all $i\ge 0$.
A step with $i$ even leaves $w$ unchanged, so $v_{i+1}\equiv w_{i+1}=w_i\equiv v_i\pmod 3$;
as $3\nmid v_i$, Lemma~\ref{lem:partner}(b) shows that the step is in case~I. The same holds for odd $i$.
Thus $s_{i+1}=As_i$ for even $i$ and $s_{i+1}=Bs_i$ for odd $i$, where now
\[
A(u,v,w)=\Bigl(\frac{u+v}{2},\ \frac{3u-v}{2},\ w\Bigr),\qquad
B(u,v,w)=\Bigl(\frac{u+w}{2},\ v,\ \frac{3u-w}{2}\Bigr).
\]
Again $A^2=B^2=\mathrm{id}$, and with $R=BA$,
\[
R(u,v,w)=\Bigl(\frac{u+v+2w}{4},\ \frac{3u-v}{2},\ \frac{3u+3v-2w}{4}\Bigr).
\]
As in the proof of Theorem~\ref{thm:A}, finiteness of $\cT(N)$ gives $r\ge 1$ with $R^rs_0=s_0$, and $R^js_0\in\Z^3$ for all $j\in\Z$.
Let $\pi(u,v,w)=(u-v,\ u-w)=(y_1,y_2)$. Then
\[
\frac{u+v+2w}{4}-\frac{3u-v}{2}=\frac{-3y_1-2y_2}{4},\qquad
\frac{u+v+2w}{4}-\frac{3u+3v-2w}{4}=\frac{2y_1-4y_2}{4},
\]
so $\pi\circ R=M_2\circ\pi$. By Lemma~\ref{lem:2adic}, $\pi(s_0)=0$, so $s_0$ is diagonal, a contradiction.
\end{proof}

\begin{remark}\label{rem:rotation}
With respect to the quadratic form in question, $A$ and $B$ are the reflections in the planes $a=v$ and $b=v$
(respectively $u=v$ and $u=w$), and $R=BA$ is a rotation about the line spanned by $(1,1,1)$.
Lemma~\ref{lem:2adic} can be replaced by the following argument: from $R^rs_0=s_0$ one gets $M^ry=y$ for $y=\pi(s_0)$;
the characteristic polynomial of $M$ is $\lambda^2-(\operatorname{tr}M)\lambda+1$ with $\operatorname{tr}M_1=1/4$ and $\operatorname{tr}M_2=-7/4$,
so the eigenvalues of $M$ are not algebraic integers, hence not roots of unity, and $M^r-1$ is invertible.
We have kept Lemma~\ref{lem:2adic} because it only uses congruences modulo $4$.
\end{remark}

\begin{example}\label{ex:walk}
Let $N=655=24\cdot 27+7$ and $s_0=(7,11,9)\in\cS(655)$, so $v_0\equiv 1\pmod 3$. The sequence of the proof is
\[
(7,11,9)\xrightarrow{\ \alpha,\ \mathrm{I}\ }(13,9,9)\xrightarrow{\ \beta,\ \mathrm{I}\ }(7,9,11)
\xrightarrow{\ \alpha,\ \mathrm{III}\ }(17,1,11).
\]
The last triple has $v=17\equiv 5\pmod 6$, and $17=6\cdot 3-1$, $1=2\cdot 0+1$, $11=2\cdot 5+1$ give
$27=\Tr(0)+\Tr(5)+p_5(3)$.
\end{example}

The next lemma deals with lower bounds on $a$.

\begin{lemma}\label{lem:large}
Let $s=(v,a,b)\in\cS(N)$ with $v>3a+6b$, and let $\alpha,\beta$ be as in the proof of Theorem~\ref{thm:A}.
Then $\beta(\alpha(s))=(v^*,a^*,b^*)\in\cS(N)$ satisfies $v^*\equiv v\pmod 3$ and $a^*\ge (v-a)/2$.
\end{lemma}

\begin{proof}
Since $v>3a$, the pair $(a,v)$ is in case~II or~III. Thus $\alpha(s)=(v',a',b)$ with $v'\equiv -v\pmod 3$,
$a'\ge (v-a)/2$ and $v'\ge (v-3a)/2>3b$.
Since $v'>3b$, the pair $(b,v')$ is in case~II or~III, so $\beta(\alpha(s))=(v^*,a',b^*)$ with $v^*\equiv -v'\equiv v\pmod 3$.
\end{proof}

The last three results of this section refine Theorems~\ref{thm:A} and~\ref{thm:B}: the residue $1$ can be reached with a value different from $1$.
For an odd positive integer $v$ with $v\equiv 2\pmod 3$ or $v=1$ put
\[
\hat v=\begin{cases} v & \text{if } v\equiv 2\pmod 3,\\ -1 & \text{if } v=1.\end{cases}
\]
We call an odd positive integer $v$ \emph{high} if $v\equiv 1\pmod 3$ and $v>1$, and \emph{low} if $v\equiv 2\pmod 3$ or $v=1$;
thus $\hat v$ is defined for the low integers, and every odd positive integer prime to $3$ is high or low.

\begin{lemma}\label{lem:hat}
Let $(u,v)\in\cP$ with $v$ low and $(u,v)\ne(1,1)$, and let $(u',v')=\varphi(u,v)$. Then either $v'$ is high, or $v'$ is low and
\[
(u',\hat v')=\Bigl(\frac{u+\hat v}{2},\ \frac{3u-\hat v}{2}\Bigr).
\]
\end{lemma}

\begin{proof}
Let first $v\equiv 2\pmod 3$, so $\hat v=v$. In case~I we have $v'=(3u-v)/2\equiv 2\pmod 3$ and $u'=(u+v)/2$, which is the displayed formula.
In case~II we have $u'=(u+v)/2$ and $v'=(v-3u)/2\equiv 1\pmod 3$; if $v'>1$ then $v'$ is high, and if $v'=1$ then $\hat v'=-1=(3u-v)/2$.
In case~III we have $v'=(3u+v)/2\equiv 1\pmod 3$ and $v'>1$, so $v'$ is high.

Let now $v=1$, so $\hat v=-1$ and $u\ge 3$. If $u\equiv 1\pmod 4$, the pair is in case~I and $v'=(3u-1)/2\equiv 1\pmod 3$ is high.
If $u\equiv 3\pmod 4$, the pair is in case~III, so $u'=(u-1)/2=(u+\hat v)/2$ and $v'=(3u+1)/2\equiv 2\pmod 3$, hence $\hat v'=v'=(3u-\hat v)/2$.
\end{proof}

\begin{theorem}\label{thm:C}
Let $N$ be a positive integer different from $7$ and $31$. If $\cS(N)$ contains a triple that is not diagonal,
then $\cS(N)$ contains a triple $(v,a,b)$ with $v\equiv 1\pmod 3$ and $v>1$.
\end{theorem}

\begin{proof}
Suppose that no triple of $\cS(N)$ has a high first entry. Let $\alpha,\beta,A,B,R,\pi$ be as in the proof of Theorem~\ref{thm:A},
and for $s=(v,a,b)\in\cS(N)$ put $\hat s=(\hat v,a,b)$.

We first show that $\cS(N)$ contains no triple of the form $(1,1,b)$ or $(1,a,1)$. By symmetry let $s=(1,1,b)\in\cS(N)$, so $N=4+3b^2$.
As $N\ne 7,31$ we have $b\ge 5$. If $b\equiv 1\pmod 4$, then $\beta(s)=\bigl((3b-1)/2,\,1,\,(b+1)/2\bigr)$ and $(3b-1)/2$ is high.
If $b=4c+3$ with $c\ge 1$, then $\beta(s)=(6c+5,\,1,\,2c+1)$, and
$\varphi(1,6c+5)$ equals $(3c+3,\,3c+1)$ for even $c$ (case~II) and $(3c+2,\,3c+4)$ for odd $c$ (case~III);
so the first entry of $\alpha(\beta(s))$ is $3c+1$ with $c\ge 2$, or $3c+4$, and it is high. This contradicts our assumption.

Now let $s_0\in\cS(N)$ be not diagonal and define $s_i$ by $s_{i+1}=\alpha(s_i)$ for even $i$ and $s_{i+1}=\beta(s_i)$ for odd $i$.
By Lemma~\ref{lem:hat}, applied to the pairs $(a_i,v_i)$ and $(b_i,v_i)$, we get $\hat s_{i+1}=A\hat s_i$ for even $i$ and $\hat s_{i+1}=B\hat s_i$ for odd $i$.
Hence $R^j\hat s_0=\hat s_{2j}$ for all $j\ge 0$. These vectors lie in the finite set of all $x\in\Z^3$ with $x_1^2+3x_2^2+3x_3^2=N$.
As in the proof of Theorem~\ref{thm:A} it follows that $R^j\hat s_0\in\Z^3$ for all $j\in\Z$ and then, by Lemma~\ref{lem:2adic}, that $\pi(\hat s_0)=0$,
that is, $\hat v_0=a_0=b_0$. In particular $\hat v_0>0$, so $\hat v_0=v_0$ and $s_0$ is diagonal, a contradiction.
\end{proof}

\begin{theorem}\label{thm:D}
Let $N$ be a positive integer different from $5$ and $29$. If $\cT(N)$ contains a triple that is not diagonal,
then $\cT(N)$ contains a triple $(u,v,w)$ such that $v\equiv 1\pmod 3$, $v>1$, and $w\equiv 2\pmod 3$ or $w=1$.
\end{theorem}

\begin{proof}
The set $\cT(N)$ is symmetric in $v$ and $w$, so we have to find a triple in which exactly one of $v,w$ is high.
Suppose there is none. Let $\alpha,\beta,A,B,R,\pi$ be as in the proof of Theorem~\ref{thm:B}, let $s_0=(u_0,v_0,w_0)\in\cT(N)$ be not diagonal,
and define $s_i=(u_i,v_i,w_i)$ by $s_{i+1}=\alpha(s_i)$ for even $i$ and $s_{i+1}=\beta(s_i)$ for odd $i$.

\emph{Case 1: $v_0$ and $w_0$ are high.} A move changes only one of the entries $v,w$, and by assumption the changed entry is again high.
So all $v_i,w_i$ are high, every step preserves the residue of the changed entry modulo $3$, and by Lemma~\ref{lem:partner}(b) every step is in case~I.
Thus $s_{i+1}=As_i$ for even $i$ and $s_{i+1}=Bs_i$ for odd $i$, and the proof of Theorem~\ref{thm:B} shows that $s_0$ is diagonal, a contradiction.

\emph{Case 2: $v_0$ and $w_0$ are low.} As before, all $v_i,w_i$ are low. Put $\hat s_i=(u_i,\hat v_i,\hat w_i)$.
We first show that no triple of $\cT(N)$ has the form $(1,1,w)$ or $(1,v,1)$. By symmetry let $s=(1,1,w)\in\cT(N)$, so $N=4+w^2$; by assumption $w$ is low.
As $N\ne 5,29$ we have $w\equiv 5\pmod 6$ and $w\ge 11$.
Then $\varphi(1,w)$ equals $\bigl((w+1)/2,\,(w-3)/2\bigr)$ if $w\equiv 1\pmod 4$ (case~II, and $w\ge 17$) and $\bigl((w-1)/2,\,(w+3)/2\bigr)$ if $w\equiv 3\pmod 4$ (case~III).
In both cases the second entry is high, so $\beta(s)$ has the low entry $v=1$ and a high entry $w$, contrary to our assumption.

Hence Lemma~\ref{lem:hat} applies to the pairs $(u_i,v_i)$ and $(u_i,w_i)$ and gives $\hat s_{i+1}=A\hat s_i$ for even $i$ and $\hat s_{i+1}=B\hat s_i$ for odd $i$.
The vectors $R^j\hat s_0=\hat s_{2j}$ lie in the finite set of all $x\in\Z^3$ with $3x_1^2+x_2^2+x_3^2=N$, and as in the proof of Theorem~\ref{thm:B} we obtain $\pi(\hat s_0)=0$,
that is, $u_0=\hat v_0=\hat w_0$. Then $\hat v_0,\hat w_0>0$, so $s_0$ is diagonal, a contradiction.
\end{proof}

\begin{remark}\label{rem:exceptions}
The exceptions are genuine: $\cS(31)=\{(1,1,3),(1,3,1),(5,1,1)\}$ and $\cT(29)=\{(1,1,5),(1,5,1),(3,1,1)\}$.
In the proofs of Theorems~\ref{thm:C} and~\ref{thm:D} the vector $\hat s$ differs from $s$ by the sign of the entries equal to $1$.
With this sign the moves that do not produce a high entry are given by the linear maps $A$ and $B$ alone, and this is what makes the periodicity argument applicable.
\end{remark}

\section{Primitive representations}\label{sec:seed}

An integer solution of a ternary equation is called \emph{primitive} if its three entries have no common prime factor.

\begin{proposition}\label{prop:classical}
\begin{enumerate}
\item[(a)] For every positive integer $N\equiv 7\pmod{24}$ the equation $v^2+3a^2+3b^2=N$ has a primitive solution $(v,a,b)\in\Z^3$.
\item[(b)] For every positive integer $N\equiv 5\pmod{24}$ the equation $3u^2+v^2+w^2=N$ has a primitive solution $(u,v,w)\in\Z^3$.
\end{enumerate}
\end{proposition}

\begin{proof}
Let $f(x_1,x_2,x_3)=c_1x_1^2+c_2x_2^2+c_3x_3^2$ with $(c_1,c_2,c_3)=(1,3,3)$ in case (a) and $(c_1,c_2,c_3)=(1,1,3)$ in case (b),
and let $L=\Z^3$ be the corresponding lattice in the quadratic space $(\Q^3,f)$.
We use two classical facts.

\emph{Fact 1.} Let $N$ be a positive integer. If for every prime $p$ there is $x_p\in\Z_p^3\setminus p\Z_p^3$ with $f(x_p)=N$,
then there are a lattice $L'$ in the genus of $L$ and a primitive vector $x\in L'$ with $f(x)=N$.
This is the primitive form of the local--global principle for representations by a genus;
the non-primitive form is \cite[p.~129]{Cassels}, see also \cite[102:5]{OMeara}. We recall the argument.
Since $f$ is positive definite and $N>0$, the form $f$ represents $N$ over $\R$, and by assumption over every $\Q_p$;
by the Hasse--Minkowski theorem there is $x\in\Q^3$ with $f(x)=N$.
By Witt's theorem there is, for each $p$, an isometry $\sigma_p$ of the quadratic space $(\Q_p^3,f)$ with $\sigma_p(x_p)=x$.
For all but finitely many $p$ the vector $x$ is itself a primitive vector of $L_p=\Z_p^3$, and for these $p$ we take $x_p=x$ and $\sigma_p=\mathrm{id}$.
There is a lattice $L'$ on $\Q^3$ with $L'_p=\sigma_p(L_p)$ for all $p$ \cite[81:14]{OMeara}. It lies in the genus of $L$,
and $x$ is a primitive vector of $L'_p$ for every $p$, hence a primitive vector of $L'$.

\emph{Fact 2.} The genus of $L$ consists of a single class.
This is classical; see the tables \cite{BI}, or \cite{LK}, where all single-class genera of positive ternary lattices are determined
(see also \cite{JKS,JP}). For the reader's convenience we indicate a direct verification.
By the reduction theory of positive ternary forms \cite[Chapter~12]{Cassels},
every positive definite ternary lattice with integral Gram matrix of determinant $D$ has a basis whose Gram matrix $(g_{ij})$ satisfies
\[
0<g_{11}\le g_{22}\le g_{33},\qquad 2|g_{12}|\le g_{11},\quad 2|g_{13}|\le g_{11},\quad 2|g_{23}|\le g_{22},\qquad g_{11}g_{22}g_{33}\le 2D.
\]
If $g_{11}=1$, then $g_{12}=g_{13}=0$ and the lattice is $\langle 1\rangle\perp K$ with $K$ a binary lattice of determinant $D$ whose Gram matrix satisfies
$2|g_{23}|\le g_{22}\le g_{33}$; for $D=3$ this gives $y^2+3z^2$ and $2y^2+2yz+2z^2$, and for $D=9$ it gives $y^2+9z^2$, $2y^2+2yz+5z^2$ and $3y^2+3z^2$.
If $g_{11}\ge 2$, then $g_{11}g_{22}g_{33}\ge 8$, which excludes $D=3$; for $D=9$ the diagonal is $(2,2,2)$, $(2,2,3)$, $(2,2,4)$ or $(2,3,3)$ with $|g_{ij}|\le 1$ for $i\ne j$, and
\[
\det(g_{ij})=g_{11}g_{22}g_{33}+2g_{12}g_{13}g_{23}-g_{11}g_{23}^2-g_{22}g_{13}^2-g_{33}g_{12}^2=9
\]
holds only for the diagonal $(2,2,3)$ with $g_{12}=\pm1$, $g_{13}=g_{23}=0$, that is, for $2x^2+2xy+2y^2+3z^2$.
So there are two classes of determinant $3$, namely $x^2+y^2+3z^2$ and $x^2+2y^2+2yz+2z^2$, and four classes of determinant $9$, namely
$x^2+y^2+9z^2$, $x^2+3y^2+3z^2$, $x^2+2y^2+2yz+5z^2$ and $2x^2+2xy+2y^2+3z^2$.
Over $\Z_3$ these six forms are equivalent to
$\langle 1,1\rangle\perp\langle 3\rangle$, $\langle 1,2\rangle\perp\langle 6\rangle$, and
$\langle 1,1\rangle\perp\langle 9\rangle$, $\langle 1\rangle\perp\langle 3,3\rangle$, $\langle 1,2\rangle\perp\langle 18\rangle$, $\langle 2\rangle\perp\langle 3,6\rangle$ respectively.
Two forms of the same determinant in this list either have Jordan components of different ranks, or unimodular components whose determinants $1$ and $2$ differ by a non-square unit;
by the uniqueness of the Jordan type over $\Z_3$ \cite[92:2]{OMeara} they are not equivalent over $\Z_3$.
Hence each of the six forms is alone in its genus.

It remains to verify the hypothesis of Fact~1.

Let $p\ge 5$. If $N\equiv c_3\pmod p$, then $N/c_3$ is the square of a unit of $\Z_p$, and $f(0,0,x_3)=N$ has a solution with $x_3$ a unit.
Otherwise the congruence $c_1x_1^2+c_2x_2^2\equiv N-c_3\pmod p$ has a solution, because each of the sets $\{c_1x_1^2\}$ and $\{N-c_3-c_2x_2^2\}$ has $(p+1)/2$ elements modulo $p$;
this solution is not $(0,0)$, so by Hensel's lemma it lifts to $x_1,x_2\in\Z_p$ with $f(x_1,x_2,1)=N$.

Let $p=3$. In case (a), $N\equiv 1\pmod 3$ is the square of a unit, so $f(x_1,0,0)=x_1^2=N$ is solvable.
In case (b), $N-1\equiv 1\pmod 3$ is the square of a unit, so $f(x_1,1,0)=x_1^2+1=N$ is solvable.

Let $p=2$. In case (a), $N-6\equiv 1\pmod 8$ is the square of a unit, so $f(x_1,1,1)=x_1^2+6=N$ is solvable.
In case (b), $N-4\equiv 1\pmod 8$ is the square of a unit, so $f(x_1,1,1)=x_1^2+4=N$ is solvable.

In all cases the solution has a unit entry. By Fact~1 there are $L'$ in the genus of $L$ and a primitive $x\in L'$ with $f(x)=N$,
and by Fact~2 there is an isometry $L'\to L$; it maps $x$ to a primitive vector of $L=\Z^3$, which is the required solution
(in case (b) with $(x_1,x_2,x_3)=(v,w,u)$).
\end{proof}

\begin{lemma}\label{lem:seed}
\begin{enumerate}
\item[(a)] If $N\equiv 7\pmod{24}$ and $N>7$, then $\cS(N)$ contains a triple that is not diagonal.
\item[(b)] If $N\equiv 5\pmod{24}$ and $N>5$, then $\cT(N)$ contains a triple that is not diagonal.
\end{enumerate}
\end{lemma}

\begin{proof}
(a) By Proposition~\ref{prop:classical}(a) there is a primitive solution of $v^2+3a^2+3b^2=N$; changing signs we may assume $v,a,b\ge 0$.
As $N\equiv 1\pmod 3$, every solution has $3\nmid v$.
Squares are congruent to $0$, $1$ or $4$ modulo $8$, and $N\equiv 7\pmod 8$.
If $v$ is odd, then $3(a^2+b^2)\equiv 6$, so $a^2+b^2\equiv 2\pmod 8$ and $a,b$ are odd.
If $v$ is even, then $3(a^2+b^2)\equiv 7-v^2$, so $a^2+b^2\equiv 5\pmod 8$ if $4\mid v$ and $a^2+b^2\equiv 1\pmod 8$ if $v\equiv 2\pmod 4$;
hence exactly one of $a,b$ is even, and the even one is $\equiv 2\pmod 4$ if $4\mid v$ and divisible by $4$ if $v\equiv 2\pmod 4$.

\emph{Case 1: $v,a,b$ are odd.} Then $(v,a,b)\in\cS(N)$. If it were diagonal, primitivity would give $v=a=b=1$ and $N=7$.

\emph{Case 2: $v$ is even.} By the symmetry of the equation in $a$ and $b$ we may assume that $b$ is even. Then $v/2+b/2$ is odd. Put
\[
b_1=\frac{|b-v|}{2},\qquad v_1=\frac{3b+v}{2}.
\]
These are odd positive integers with $3b_1^2+v_1^2=3b^2+v^2$, so $(v_1,a,b_1)$ is a solution in odd positive integers and belongs to $\cS(N)$.
Suppose $v_1=a=b_1=k$. Then $3b+v=2k$ and $|b-v|=2k$ with $b,v\ge 0$, which forces $b=0$ and $v=2k$.
So $(v,a,b)=(2k,k,0)$, and primitivity gives $k=1$ and $N=7$.

(b) By Proposition~\ref{prop:classical}(b) there is a primitive solution of $3u^2+v^2+w^2=N$ with $u,v,w\ge 0$.
As $N\equiv 2\pmod 3$, every solution has $3\nmid vw$. Here $N\equiv 5\pmod 8$.
If $u$ is odd, then $v^2+w^2\equiv 2\pmod 8$ and $v,w$ are odd.
If $u$ is even, then $v^2+w^2\equiv 5\pmod 8$ if $4\mid u$ and $v^2+w^2\equiv 1\pmod 8$ if $u\equiv 2\pmod 4$;
hence exactly one of $v,w$ is even, and $u/2$ plus half of the even one is odd.

\emph{Case 1: $u,v,w$ are odd.} Then $(u,v,w)\in\cT(N)$, and it is not diagonal since it is primitive and $N>5$.

\emph{Case 2: $u$ is even.} We may assume that $v$ is even. Put $u_1=|u-v|/2$ and $v_1=(3u+v)/2$.
These are odd positive integers with $3u_1^2+v_1^2=3u^2+v^2$, so $(u_1,v_1,w)\in\cT(N)$.
If $u_1=v_1=w=k$, then $3u+v=2k$ and $|u-v|=2k$ force $u=0$, $v=2k$, so $(u,v,w)=(0,2k,k)$ and $k=1$, $N=5$.
\end{proof}

\begin{remark}
The passage from an even solution to an odd one in Case~2 is the move of \cite[Lemma~3.2]{S15}.
Lemma~\ref{lem:seed} is the only place where primitivity is used, and it is needed only for $N=7k^2$ in (a) and $N=5k^2$ in (b):
for such $N$ the solutions $k\cdot(\pm1,\pm1,\pm1)$, $k\cdot(\pm2,\pm1,0)$, $k\cdot(\pm 2,0,\pm 1)$
(respectively $k\cdot(\pm1,\pm1,\pm1)$, $k\cdot(0,\pm2,\pm1)$, $k\cdot(0,\pm 1,\pm 2)$) lead to the diagonal triple $(k,k,k)$, which the maps $\alpha$ and $\beta$ fix.
For all other $N$ the proof of Lemma~\ref{lem:seed} works with any integer solution, primitive or not.
\end{remark}

\section{Proofs of the theorems}\label{sec:proofs}

For $k\ge 1$, $l>-k$, $k\equiv l\pmod 2$ and $x\in\N$ we have
\begin{equation}\label{eq:square}
8k\,\psi_{k,l}(x)+l^2=(2kx+l)^2 .
\end{equation}

\begin{lemma}\label{lem:reduce}
Let $e\in\{-1,1,5\}$, let $d$ be a positive odd integer and $n\in\N$. Put
\[
N=24n+e^2+3d^2+3 .
\]
Then $N\equiv 7\pmod{24}$, and $n=\psi_{3,e}(x)+\psi_{1,d}(y)+\Tr(z)$ for some $x,y,z\in\N$
if and only if there is $(v,a,b)\in\cS(N)$ with $\max\{a,b\}\ge d$ and
\[
\begin{cases}
v\equiv 2\pmod 3 \ \text{ or } \ v=1 & \text{if } e=-1,\\
v\equiv 1\pmod 3 & \text{if } e=1,\\
v\equiv 2\pmod 3 & \text{if } e=5.
\end{cases}
\]
\end{lemma}

\begin{proof}
As $d$ is odd, $3d^2\equiv 3\pmod{24}$, and $e^2\equiv 1\pmod{24}$; so $N\equiv 7\pmod{24}$.
By \eqref{eq:square},
\[
24\bigl(\psi_{3,e}(x)+\psi_{1,d}(y)+\Tr(z)\bigr)+e^2+3d^2+3=(6x+e)^2+3(2y+d)^2+3(2z+1)^2 .
\]
If $n$ is represented, put $v=|6x+e|$, $a=2y+d$, $b=2z+1$. Then $(v,a,b)\in\cS(N)$, $a\ge d$, and
$v=6x+1$ for $e=1$, $v=6x+5$ for $e=5$, while for $e=-1$ either $x=0$ and $v=1$, or $x\ge 1$ and $v=6x-1$.
Conversely, let $(v,a,b)\in\cS(N)$ be as in the statement. The set $\cS(N)$ is symmetric in $a$ and $b$, so we may assume $a\ge d$.
Put $y=(a-d)/2$ and $z=(b-1)/2$. As $v$ is odd, $v\equiv 1\pmod 3$ means $v\equiv 1\pmod 6$ and $v\equiv 2\pmod 3$ means $v\equiv 5\pmod 6$.
Put $x=(v-1)/6$ if $e=1$; $x=(v-5)/6$ if $e=5$; and for $e=-1$ put $x=(v+1)/6$ if $v\equiv 5\pmod 6$ and $x=0$ if $v=1$.
Then $x,y,z\in\N$, $(6x+e)^2=v^2$, and the displayed identity gives $n=\psi_{3,e}(x)+\psi_{1,d}(y)+\Tr(z)$.
\end{proof}

\begin{proposition}\label{prop:family}
Let $e\in\{-1,1,5\}$, let $d$ be a positive odd integer, $n\in\N$ and $N=24n+e^2+3d^2+3$.
If $d=1$, or if $N>87(d-2)^2$, then $n=\psi_{3,e}(x)+\psi_{1,d}(y)+\Tr(z)$ for some $x,y,z\in\N$.
\end{proposition}

\begin{proof}
If $N=7$, then $e=\pm1$, $d=1$ and $n=0$, and $0=\psi_{3,e}(0)+\psi_{1,1}(0)+\Tr(0)$. Let $N>7$.
Let $\rho=1$ if $e=1$ and $\rho=2$ if $e\in\{-1,5\}$.
By Lemma~\ref{lem:seed}(a) and Theorem~\ref{thm:A} there is $s=(v,a,b)\in\cS(N)$ with $v\equiv \rho\pmod 3$.
If $\max\{a,b\}\ge d$, we are done by Lemma~\ref{lem:reduce}; this is always the case when $d=1$.
Otherwise $d\ge 3$ and $a,b\le K:=d-2$. If $N>87K^2$, then
\[
v^2=N-3a^2-3b^2>87K^2-6K^2=81K^2 ,
\]
so $v>9K\ge 3a+6b$. By Lemma~\ref{lem:large} there is $(v^*,a^*,b^*)\in\cS(N)$ with $v^*\equiv v\equiv \rho\pmod 3$ and
$a^*\ge (v-a)/2>(9K-K)/2=4K\ge K+2=d$, and Lemma~\ref{lem:reduce} applies again.
\end{proof}

\begin{proof}[Proof of Theorem~\ref{thm:family}]
\emph{Universality.} Let $(e,d)$ be one of the pairs with $e=-1$, $d\le 17$; $e=1$, $d\le 7$; $e=5$, $d\le 3$.
By Proposition~\ref{prop:family} every $n$ with $24n+e^2+3d^2+3>87(d-2)^2$ is represented, and for $d=1$ every $n$ is.
For $d\ge 3$ the remaining $n$ satisfy
\[
n\le \frac{87(d-2)^2-e^2-3d^2-3}{24}.
\]
There are twelve such pairs $(e,d)$, the bound is largest for $(e,d)=(-1,17)$, where it is $779$, and altogether $2468$ triples $(e,d,n)$ have to be treated.
For each of them a representation $n=\psi_{3,e}(x)+\psi_{1,d}(y)+\Tr(z)$ with $x,y,z\in\N$ exists; it is found at once by trying all $x$ and $y$,
and the list of these representations is part of the files accompanying this paper (see Section~\ref{sec:remarks}).

\emph{Non-universality.} This is contained in \cite[Theorem~1.4(ii)]{S20}; here is a direct check.
If $y\ge 1$ then $\psi_{1,d}(y)\ge (d+1)/2$.
Let $e=-1$ and $d\ge 19$. Then $\psi_{1,d}(y)\ge 10$ for $y\ge 1$, so a representation of $9$ has $y=0$, and
$9-p_5(x)\in\{9,8,4\}$ is not triangular. Hence $9$ is not represented.
Let $e=1$ and $d\ge 9$. Then $\psi_{1,d}(y)\ge 5$ for $y\ge 1$, and $4-\pb(x)\in\{4,2\}$ is not triangular; so $4$ is not represented.
Let $e=5$ and $d\ge 5$. Then $\psi_{1,d}(y)\ge 3$ for $y\ge 1$ and $\psi_{3,5}(x)\ge 4$ for $x\ge 1$, and $2$ is not triangular; so $2$ is not represented.
\end{proof}

\begin{proof}[Proof of Theorem~\ref{thm:main}]
The first assertion is Theorem~\ref{thm:family}(i) for $d=1$.
For the second, let $n\ge 1$. By Theorem~\ref{thm:family}(iii) for $d=1$ there are $x,y,z\in\N$ with
$n-1=\psi_{3,5}(z)+\Tr(x)+\Tr(y)$, and $\psi_{3,5}(z)+1=(z+1)(3z+2)/2=p_5(z+1)$.
\end{proof}

\begin{lemma}\label{lem:TT}
$\{\Tr(y)+\Tr(z):y,z\in\N\}=\{y(y+1)+z^2:y,z\in\N\}$.
\end{lemma}

\begin{proof}
This is well known, see \cite[(1.5)]{S15}: if $y\ge z$ then
$\Tr(y)+\Tr(z)=r(r+1)+s^2$ with $\{2r+1,2s\}=\{y+z+1,y-z\}$, and conversely
$r(r+1)+s^2=\Tr(r+s)+\Tr(|r-s|-[r<s])$, where $[r<s]$ is $1$ if $r<s$ and $0$ otherwise.
\end{proof}

\begin{proof}[Proof of Corollary~\ref{cor:square}]
By Theorem~\ref{thm:family} for $d=1$ and Lemma~\ref{lem:TT}, for $e\in\{-1,1,5\}$ every $n\in\N$ is of the form
$\psi_{3,e}(x)+y(y+1)+z^2=\psi_{3,e}(x)+\psi_{2,2}(y)+\psi_{2,0}(z)$. For $e=-1$ this is $p_5(x)+2p_3(y)+p_4(z)$.
\end{proof}

\begin{proof}[Proof of Theorem~\ref{thm:mixed}]
Let $n\in\N$ and $N=24n+5$. By \eqref{eq:square},
\[
24\bigl(\Tr(x)+p_5(y)+\pb(z)\bigr)+5=3(2x+1)^2+(6y-1)^2+(6z+1)^2 .
\]
If $n=0$ take $x=y=z=0$. Let $n\ge 1$, so $N>5$.
By Lemma~\ref{lem:seed}(b) and Theorem~\ref{thm:B} there is $(u,v,w)\in\cT(N)$ with $v\not\equiv w\pmod 3$.
The set $\cT(N)$ is symmetric in $v$ and $w$, so we may assume $v\equiv 2$ and $w\equiv 1\pmod 3$, that is,
$v\equiv 5$ and $w\equiv 1\pmod 6$. Then $x=(u-1)/2$, $y=(v+1)/6\ge 1$, $z=(w-1)/6$ are in $\N$ and
$n=\Tr(x)+p_5(y)+\pb(z)$. This proves the first two assertions, and hence the universality of $(3,1,3,-1,1,1)$.
Finally, for $m\in\N$ apply this to $n=m+1$: since $p_5(y)-1=(y-1)(3y+2)/2=\psi_{3,5}(y-1)$, we get
$m=\psi_{3,5}(y-1)+\pb(z)+\Tr(x)$ with $y-1,z,x\in\N$. So $(3,5,3,1,1,1)$ is universal over $\N$.
\end{proof}

\begin{proof}[Proof of Theorem~\ref{thm:shift}]
(i) Let $n\ge 2$ and $N=24n+7$. Then $N\ge 55$, and by Lemma~\ref{lem:seed}(a) the set $\cS(N)$ contains a triple that is not diagonal.
By Theorem~\ref{thm:C} there is $(v,a,b)\in\cS(N)$ with $v\equiv 1\pmod 3$ and $v>1$. As $v$ is odd, $v\equiv 1\pmod 6$ and $v\ge 7$.
Put $z=(v-1)/6\ge 1$, $x=(a-1)/2$ and $y=(b-1)/2$. By \eqref{eq:square},
$24\bigl(\Tr(x)+\Tr(y)+\pb(z)\bigr)+7=3a^2+3b^2+v^2=N$, so $n=\Tr(x)+\Tr(y)+\pb(z)$.

(ii) Let $n\ge 2$ and $N=24n+5$. Then $N\ge 53$. By Lemma~\ref{lem:seed}(b) and Theorem~\ref{thm:D} there is $(u,v,w)\in\cT(N)$
with $v\equiv 1\pmod 3$, $v>1$, and $w\equiv 2\pmod 3$ or $w=1$. Put $x=(u-1)/2$ and $z=(v-1)/6\ge 1$; put $y=(w+1)/6$ if $w\equiv 2\pmod 3$, and $y=0$ if $w=1$.
In both cases $(6y-1)^2=w^2$, and the identity in the proof of Theorem~\ref{thm:mixed} gives $n=\Tr(x)+p_5(y)+\pb(z)$.

For the tuples, note that $\pb(z)-2=(z-1)(3z+4)/2=\psi_{3,7}(z-1)$. Let $m\in\N$. Applying (i) to $n=m+2$ we get
$m=\psi_{3,7}(z-1)+\Tr(x)+\Tr(y)$, so $(3,7,1,1,1,1)$ is universal over $\N$, and by Lemma~\ref{lem:TT} so is $(3,7,2,2,2,0)$.
Applying (ii) to $n=m+2$ we get $m=\psi_{3,7}(z-1)+p_5(y)+\Tr(x)$, so $(3,7,3,-1,1,1)$ is universal over $\N$.
\end{proof}

\section{Two triangular numbers with different indices}\label{sec:distinct}

In this section we prove Theorem~\ref{thm:distinct}. We use the notation of Section~\ref{sec:descent}:
the maps $\alpha$, $\beta$, $A$, $B$ and $\pi$ from the proof of Theorem~\ref{thm:A}, and the words \emph{high} and \emph{low} introduced before Lemma~\ref{lem:hat};
a triple is called high or low if its first entry is.
Let $\sigma(v,a,b)=(v,b,a)$; then $\sigma\alpha=\beta\sigma$ and $\sigma A\sigma=B$.
A triple $(v,a,b)\in\cS(N)$ is called \emph{primitive} if $\gcd(v,a,b)=1$, \emph{symmetric} if $a=b$, and \emph{good} if $v$ is low and $a\ne b$.

\begin{theorem}\label{thm:E}
Let $N$ be a positive integer different from $7$. If $\cS(N)$ contains a primitive triple, then $\cS(N)$ contains a good triple.
\end{theorem}

The proof needs two lemmas.

\begin{lemma}\label{lem:gcd}
The maps $\alpha$ and $\beta$ send primitive triples to primitive triples.
If $N\equiv 7\pmod{24}$, then $\cS(N)$ contains a primitive triple.
\end{lemma}

\begin{proof}
Let $(u,v)\in\cP$ and $(u',v')=\varphi(u,v)$. By Definition~\ref{def:partner} the numbers $2u'$ and $2v'$ are integral linear combinations of $u$ and $v$,
and by Lemma~\ref{lem:partner}(c) the numbers $2u$ and $2v$ are integral linear combinations of $u'$ and $v'$. As all four numbers are odd,
an integer divides $u$ and $v$ if and only if it divides $u'$ and $v'$. This proves the first assertion.
For the second, consider the triple constructed in the proof of Lemma~\ref{lem:seed}(a) from a primitive solution $(v,a,b)$.
In Case~1 it is $(v,a,b)$ itself. In Case~2 it is $(v_1,a,b_1)$ with $b_1=|b-v|/2$ and $v_1=(3b+v)/2$; here $2b$ and $2v$ are integral linear combinations of $b_1$ and $v_1$,
so an odd common divisor of $v_1,a,b_1$ divides $v,a,b$ and is $1$.
\end{proof}

\begin{lemma}\label{lem:val}
Let $x_0\in\Q^3$ with $\pi(x_0)=(y,y)$, $y\ne 0$, and let $x_{j+1}=Ax_j$ for even $j$ and $x_{j+1}=Bx_j$ for odd $j$.
Then for every $j\ge 1$ the two entries of $\pi(x_j)$ are nonzero and have different $2$-adic valuations. In particular $x_j$ is not symmetric.
\end{lemma}

\begin{proof}
For $Y=(Y_1,Y_2)=\pi(x)$ one computes $\pi(Ax)=(-Y_1,\ Y_2-\tfrac32Y_1)$ and $\pi(Bx)=(Y_1-\tfrac32Y_2,\ -Y_2)$.
Let $e$ be the $2$-adic valuation of $y$. We claim that the valuations of the entries of $\pi(x_j)$ are $(e-j+1,\ e-j)$ for odd $j$ and $(e-j,\ e-j+1)$ for even $j\ge 2$.
For $j=1$ we have $\pi(x_1)=(-y,-y/2)$. If $j$ is odd and the claim holds for $j$, then $\tfrac32Y_2$ has valuation $e-j-1<e-j+1$,
so $\pi(x_{j+1})=\pi(Bx_j)$ has valuations $(e-j-1,\ e-j)$. If $j\ge 2$ is even and the claim holds for $j$, then $\tfrac32Y_1$ has valuation $e-j-1<e-j+1$,
so $\pi(x_{j+1})=\pi(Ax_j)$ has valuations $(e-j,\ e-j-1)$.
\end{proof}

\begin{proof}[Proof of Theorem~\ref{thm:E}]
Let $t\in\cS(N)$ be primitive, and let $O\subseteq\cS(N)$ be the set of all triples obtained from $t$ by applying $\alpha$ and $\beta$ repeatedly.
By Lemma~\ref{lem:gcd} all triples in $O$ are primitive; as $N\ne 7$, none of them is diagonal. By Lemma~\ref{lem:partner}(c), $\alpha$ and $\beta$ are involutions of the finite set $O$.
Suppose that $O$ contains no good triple. Then every triple in $O$ is symmetric or has a high first entry.

\emph{Step 1: symmetric triples.} Let $s=(v,a,a)\in O$ with $v\equiv 2\pmod 3$. Then $v\ne a$.
The pair $(a,v)$ is not in case~I, because otherwise $\alpha(s)=\bigl((3a-v)/2,(a+v)/2,a\bigr)$ would be good.
We say that $s$ is of \emph{type} II if $(a,v)$ is in case~II, of type $\mathrm{III}^+$ if it is in case~III and $a>v$, and of type $\mathrm{III}^-$ if it is in case~III and $a<v$. Accordingly
\[
\alpha(s)=\Bigl(\frac{v-3a}{2},\frac{a+v}{2},a\Bigr),\qquad \Bigl(\frac{3a+v}{2},\frac{a-v}{2},a\Bigr),\qquad \Bigl(\frac{3a+v}{2},\frac{v-a}{2},a\Bigr).
\]
The first entry of $\alpha(s)$ is $\equiv 1\pmod 3$, and $\alpha(s)$ is not symmetric, since $v\ne a$, $v>0$ and $3\nmid v$; so $\alpha(s)$ is high, and so is $\beta(s)=\sigma\alpha(s)$.
Put $s^-=(-v,a,a)$. If $s$ is of type $\mathrm{III}^+$, then
\begin{equation}\label{eq:minus}
\alpha(s)=As^-\qquad\text{and}\qquad \beta(s)=Bs^- .
\end{equation}

\emph{Step 2: a run.} By the proof of Theorem~\ref{thm:A}, applied to $t$, the set $O$ contains a triple whose first entry is $\equiv 2\pmod 3$; it is low, hence symmetric.
Call it $s_0$ and define $s_{i+1}=\alpha(s_i)$ for even $i$ and $s_{i+1}=\beta(s_i)$ for odd $i$. The map sending $(s,0)$ to $(\alpha(s),1)$ and $(s,1)$ to $(\beta(s),0)$ is a bijection of the finite set $O\times\{0,1\}$,
and $(s_i,\,i\bmod 2)$ is the orbit of $(s_0,0)$ under it. Hence $s_{2m}=s_0$ for some $m\ge 1$; in particular $s_i$ is symmetric for some $i\ge 1$. By Step~1, $s_1$ is not symmetric. Let $k\ge 1$ be such that $s_1,\dots,s_k$ are not symmetric and $s':=s_{k+1}$ is symmetric.
Then $s_1,\dots,s_k$ are high. For $1\le i<k$ the step from $s_i$ to $s_{i+1}$ preserves the residue of the first entry, so it is in case~I by Lemma~\ref{lem:partner}(b):
$s_{i+1}=Bs_i$ for odd $i$ and $s_{i+1}=As_i$ for even $i$. Let $X=B$ if $k$ is odd and $X=A$ if $k$ is even.
If the first entry of $s'$ is $\equiv 1\pmod 3$, the last step is in case~I as well and $s'=Xs_k$.
Otherwise $s'$ is one of the triples of Step~1, and $s_k$ is $\alpha(s')$ or $\beta(s')$.

\emph{Step 3: one of $s_0$, $s'$ is a triple as in Step~1 of type \textup{II} or $\mathrm{III}^-$.}
Suppose not. Then $s_0$ is of type $\mathrm{III}^+$, and either $s'=Xs_k$, or $s'$ is of type $\mathrm{III}^+$ and $s_k=Xs'^-$ by \eqref{eq:minus}, that is, $Xs_k=s'^-$.
In both cases $Xs_k$ is symmetric. Put $x_0=s_0^-$, $x_i=s_i$ for $1\le i\le k$ and $x_{k+1}=Xs_k$. By \eqref{eq:minus}, $x_1=Ax_0$, so $x_{j+1}=Ax_j$ for even $j$ and $x_{j+1}=Bx_j$ for odd $j$,
and $\pi(x_0)=(a+v,a+v)\ne 0$ if $s_0=(v,a,a)$. This contradicts Lemma~\ref{lem:val}.

\emph{Step 4: reduction to a triple of type \textup{II} or $\mathrm{III}^-$.}
If $s_0$ is not of type II or $\mathrm{III}^-$, we replace the sequence $s_0,s_1,\dots,s_k,s'$ by the reversed sequence $s',s_k,\dots,s_1,s_0$,
and, if $s_k=\beta(s')$, we apply $\sigma$ to all its members. By Step~3 we obtain in all cases the following situation:
there are a symmetric triple $s=(v,a,a)\in O$ with $v\equiv 2\pmod 3$ of type II or $\mathrm{III}^-$, an integer $k\ge 1$, and vectors $x_1,\dots,x_k$ with positive entries such that
$x_1=\alpha(s)$, $x_{i+1}=Bx_i$ for odd $i<k$ and $x_{i+1}=Ax_i$ for even $i<k$, and such that the following holds.
Write $x_k=(v_k,a_k,b_k)$ and put $(p,q)=(b_k,a_k)$ for odd $k$, $(p,q)=(a_k,b_k)$ for even $k$. Then
\begin{equation}\label{eq:end}
\frac{p+v_k}{2}=q\qquad\text{or}\qquad \frac{p-v_k}{2}=q\qquad\text{or}\qquad \frac{v_k-p}{2}=q .
\end{equation}
To see this in the reversed case, put $t_i=s_{k+1-i}$. As $A^2=B^2=\mathrm{id}$, we have $t_{i+1}=Bt_i$ if $k-i$ is odd and $t_{i+1}=At_i$ if $k-i$ is even $(1\le i<k)$.
For even $k$ this is the required rule, $t_1=s_k=\alpha(s')$, and $s_0=\alpha(t_k)$. For odd $k$ the roles of $A$ and $B$ are exchanged and $t_1=\beta(s')$;
since $\sigma A\sigma=B$ and $\sigma\beta=\alpha\sigma$, the vectors $x_i=\sigma t_i$ satisfy the required rule, $x_1=\alpha(s')$, and $s_0=\sigma(s_0)=\beta(x_k)$.
In all cases the member following $x_k$ is symmetric, and it is obtained from $x_k$ by replacing the pair $(p,v_k)$ by its partner, whose first entry is $(p+v_k)/2$ or $|p-v_k|/2$. This gives \eqref{eq:end}.

\emph{Step 5: the triple $s$ does not exist.} Extend the sequence $x_1,\dots,x_k$ by the same rule to vectors $x_1=\alpha(s)$, $x_2=Bx_1$, $x_3=Ax_2$, $x_4=Bx_3$, $x_5=Ax_4$ in $\Q^3$. By the formulas of Step~1 a direct computation gives, for $s$ of type II,
\[
x_1=\Bigl(\frac{v-3a}{2},\frac{v+a}{2},a\Bigr),\quad
x_2=\Bigl(\frac{9a-v}{4},\frac{v+a}{2},\frac{v-a}{4}\Bigr),\quad
x_3=\Bigl(\frac{7v-3a}{8},\frac{v+11a}{8},\frac{v-a}{4}\Bigr),
\]
and the first entry of $x_4$ is $-(v+3a)/16<0$. For $s$ of type $\mathrm{III}^-$ it gives
\begin{gather*}
x_1=\Bigl(\frac{v+3a}{2},\frac{v-a}{2},a\Bigr),\qquad
x_2=\Bigl(\frac{3a-v}{4},\frac{v-a}{2},\frac{v+5a}{4}\Bigr),\\
x_3=\Bigl(\frac{7v-9a}{8},\frac{v+a}{8},\frac{v+5a}{4}\Bigr),\qquad
x_4=\Bigl(\frac{39a-v}{16},\frac{v+a}{8},\frac{9v+a}{16}\Bigr),
\end{gather*}
and the first entry of $x_5$ is $(7v-33a)/32$, which is negative when the first entry of $x_2$ is positive.
As $x_1,\dots,x_k$ have positive entries, we get $k\le 3$ for type II, with $v<9a$ if $k\ge 2$; and $k\le 4$ for type $\mathrm{III}^-$, with $v<3a$ if $k\ge 2$.
Recall that $v>3a$ for type II and $v>a$ for type $\mathrm{III}^-$.

\begin{table}[ht]
\small\setlength{\tabcolsep}{5pt}
\begin{tabular}{cccccc}
type & $k$ & $(v_k,p,q)$ & $\tfrac{p+v_k}{2}=q$ & $\tfrac{p-v_k}{2}=q$ & $\tfrac{v_k-p}{2}=q$ \\[2pt] \hline
II & $1$ & $\bigl(\tfrac{v-3a}{2},\,a,\,\tfrac{v+a}{2}\bigr)$ & $v=-3a$ & $v=a$ & $v=-7a$ \rule{0pt}{14pt}\\[3pt]
II & $2$ & $\bigl(\tfrac{9a-v}{4},\,\tfrac{v+a}{2},\,\tfrac{v-a}{4}\bigr)$ & $v=13a$ & $\mathbf{v=5a}$ & $5v=9a$ \\[3pt]
II & $3$ & $\bigl(\tfrac{7v-3a}{8},\,\tfrac{v-a}{4},\,\tfrac{v+11a}{8}\bigr)$ & $\mathbf{7v=27a}$ & $v=-3a$ & $\mathbf{3v=23a}$ \\[3pt]
$\mathrm{III}^-$ & $1$ & $\bigl(\tfrac{v+3a}{2},\,a,\,\tfrac{v-a}{2}\bigr)$ & $\mathbf{v=7a}$ & $3v=a$ & $\mathbf{v=3a}$ \\[3pt]
$\mathrm{III}^-$ & $2$ & $\bigl(\tfrac{3a-v}{4},\,\tfrac{v-a}{2},\,\tfrac{v+5a}{4}\bigr)$ & $v=-9a$ & $v=15a$ & $v=-a$ \\[3pt]
$\mathrm{III}^-$ & $3$ & $\bigl(\tfrac{7v-9a}{8},\,\tfrac{v+5a}{4},\,\tfrac{v+a}{8}\bigr)$ & $7v=a$ & $\mathbf{7v=17a}$ & $v=7a$ \\[3pt]
$\mathrm{III}^-$ & $4$ & $\bigl(\tfrac{39a-v}{16},\,\tfrac{v+a}{8},\,\tfrac{9v+a}{16}\bigr)$ & $\mathbf{17v=39a}$ & $5v=-13a$ & $\mathbf{3v=5a}$
\end{tabular}
\vspace{4pt}
\caption{The solutions of the equations \eqref{eq:end}. Those compatible with $v,a>0$ and with the inequalities for the type and for $k$ are in bold.}\label{tab:eq}
\end{table}
For each of these seven pairs (type, $k$) the three equations in \eqref{eq:end} are linear in $v$ and $a$. Their solutions are listed in Table~\ref{tab:eq}.
Only eight of the solutions are compatible with $v,a>0$, with $v>3a$ (type II) or $v>a$ (type $\mathrm{III}^-$), and with $v<9a$ (type II, $k\ge 2$) or $v<3a$ (type $\mathrm{III}^-$, $k\ge 2$), namely
\[
v=5a,\quad 7v=27a,\quad 3v=23a,\quad v=7a,\quad v=3a,\quad 7v=17a,\quad 17v=39a,\quad 3v=5a .
\]
As $\gcd(v,a)=1$ and $v\equiv 2\pmod 3$, the equations $7v=27a$, $v=3a$, $17v=39a$ (which give $3\mid v$) and $v=7a$ (which gives $s=(7,1,1)$) are impossible,
and $s$ is one of the triples $(5,1,1)$, $(23,3,3)$, $(17,7,7)$, $(5,3,3)$. But
\begin{align*}
&\alpha(5,1,1)=(1,3,1), &&\beta\alpha(23,3,3)=\beta(7,13,3)=(1,13,5),\\
&\beta\alpha(17,7,7)=\beta(19,5,7)=(1,5,13), &&\beta\alpha(5,3,3)=\beta(7,1,3)=(1,1,5)
\end{align*}
are good triples in $O$. This contradiction proves the theorem.
\end{proof}

\begin{remark}\label{rem:primitive}
Primitivity is needed in Theorem~\ref{thm:E}: for $N=55k^2$ with $k\equiv 5\pmod 6$ the triples $k\cdot(7,1,1)$, $k\cdot(5,3,1)$, $k\cdot(1,3,3)$, $k\cdot(5,1,3)$
are permuted by $\alpha$ and $\beta$, and none of them is good. For $k=1$ the same four triples show that the analogue of Theorem~\ref{thm:E} for the other residue fails for $N=55$;
indeed $1$ is not of the form $x(3x+1)/2+y(y+2)+z(z+1)$, and the tuple $(3,1,2,4,2,2)$ is not in the list of \cite{S20}.
Similarly, for $g\equiv 1\pmod 6$, $g>1$, the multiples by $g$ of the triples obtained from $(5,1,1)$, $(5,3,3)$ or $(17,7,7)$ by applying $\alpha$ and $\beta$ contain no good triple.
\end{remark}

\begin{proof}[Proof of Theorem~\ref{thm:distinct}]
Let $n\ge 1$ and $N=24n+7$. By Proposition~\ref{prop:classical}(a) and Lemma~\ref{lem:gcd} the set $\cS(N)$ contains a primitive triple, and $N\ne 7$.
By Theorem~\ref{thm:E} there is $(v,a,b)\in\cS(N)$ with $a\ne b$ and $v\equiv 2\pmod 3$ or $v=1$.
Put $r=(a-1)/2$, $s=(b-1)/2$, and $x=(v+1)/6$ if $v\equiv 2\pmod 3$, $x=0$ if $v=1$. Then $r\ne s$, and $n=p_5(x)+\Tr(r)+\Tr(s)$ by the identity in the proof of Lemma~\ref{lem:reduce} (with $e=-1$ and $d=1$).

For the tuple, let $m\in\N$ and write $m+1=p_5(x)+\Tr(r)+\Tr(s)$ with $r>s$. By the proof of Lemma~\ref{lem:TT}, $\Tr(r)+\Tr(s)=z(z+1)+w^2$ with $\{2z+1,2w\}=\{r+s+1,r-s\}$; in particular $w\ge 1$.
Hence $m=p_5(x)+z(z+1)+(w-1)(w+1)=\psi_{3,-1}(x)+\psi_{2,2}(z)+\psi_{2,4}(w-1)$.
\end{proof}

\section{Formalization}\label{sec:lean}

The arguments of Sections~\ref{sec:partner}, \ref{sec:descent} and~\ref{sec:proofs} and of Lemma~\ref{lem:seed}
have been formalized in Lean~4 with Mathlib (Section~\ref{sec:distinct}, that is, Theorem~\ref{thm:distinct}, has not been formalized); the development is a single file \texttt{TriPent.lean} in the accompanying files \cite{code}.
Proposition~\ref{prop:classical} is not formalized, since the theory of genera is not available in Mathlib.
It enters as two explicit hypotheses,
\begin{align*}
\texttt{PrimRep133}:\quad &\forall N\in\Z,\ N>0,\ N\equiv 7 \ (\mathrm{mod}\ 24)\ \Longrightarrow\\
&\exists\, v,a,b\in\Z:\ v^2+3a^2+3b^2=N,\ \gcd(v,a,b)=1,\\
\texttt{PrimRep113}:\quad &\forall N\in\Z,\ N>0,\ N\equiv 5 \ (\mathrm{mod}\ 24)\ \Longrightarrow\\
&\exists\, u,v,w\in\Z:\ 3u^2+v^2+w^2=N,\ \gcd(u,v,w)=1,
\end{align*}
which are exactly the two parts of Proposition~\ref{prop:classical}. The formal statements proved are the following.
\begin{itemize}
\item \texttt{exists\_other\_residue} and \texttt{exists\_mixed}: Theorems~\ref{thm:A} and~\ref{thm:B}, without any hypothesis.
\item \texttt{exists\_residue1\_gt} and \texttt{exists\_high\_low}: Theorems~\ref{thm:C} and~\ref{thm:D}, without any hypothesis.
\item \texttt{seedS} and \texttt{seedT}: Lemma~\ref{lem:seed}, starting from a primitive solution.
\item \texttt{two\_triangular\_pentagonal}: \texttt{PrimRep133} implies Theorem~\ref{thm:main}, in the form
$\forall n\in\N\ \exists x,y,z\in\N:\ n=x(x+1)/2+y(y+1)/2+z(3z-1)/2$ and ($n\ge 1\Rightarrow z\ge 1$), with the division and subtraction of natural numbers.
\item \texttt{family\_m1}, \texttt{family\_p1}, \texttt{family\_p5}: \texttt{PrimRep133} implies the universality statements of Theorem~\ref{thm:family};
the $2468$ small cases are verified by the Lean kernel from an explicit list of representations.
\item \texttt{tuple\_3\_m1\_2\_2\_2\_0} and its two analogues: Corollary~\ref{cor:square}.
\item \texttt{triangular\_two\_pentagonal\_mixed} and \texttt{tuple\_3\_5\_3\_1\_1\_1}: \texttt{PrimRep113} implies Theorem~\ref{thm:mixed}.
\item \texttt{tuple\_3\_7\_1\_1\_1\_1} and \texttt{tuple\_3\_7\_2\_2\_2\_0} (from \texttt{PrimRep133}) and \texttt{tuple\_3\_7\_3\_m1\_1\_1} (from \texttt{PrimRep113}):
the three tuples of Theorem~\ref{thm:shift}.
\end{itemize}
No axioms are used beyond the three standard ones of Lean (\texttt{propext}, \texttt{Classical.choice}, \texttt{Quot.sound}), and no \texttt{sorry} or \texttt{native\_decide}.
The periodicity argument is formalized as in the proofs above: the orbit of the map $R$ stays in a finite box, the pigeonhole principle gives a period,
and Lemma~\ref{lem:2adic} is proved by strong induction on a bound for the entries.

\section{Remarks}\label{sec:remarks}

\subsection*{Computations}
The only computation used in the proofs is the direct search in the proof of Theorem~\ref{thm:family}:
$2468$ cases with $n\le 779$, which take a fraction of a second. The program and the resulting list of representations are in the accompanying files \cite{code}.
The formulas and Table~\ref{tab:eq} in Step~5 of the proof of Theorem~\ref{thm:E} were also checked by a program,
and Theorem~\ref{thm:E} was checked for every primitive triple in $\cS(24n+7)$ with $n<1.5\cdot 10^5$ (about $9.8\cdot 10^7$ triples).
Independently of the proofs, the statements of Theorems~\ref{thm:A}, \ref{thm:B}, \ref{thm:C} and~\ref{thm:D} were checked for every triple in
$\cS(24n+7)$ and $\cT(24n+5)$ with $n<3\cdot 10^5$ (about $2.8\cdot 10^8$ and $3.2\cdot 10^8$ triples);
the sequence $s_0,s_1,\dots$ of the proofs reached the required triple after at most $19$ steps,
or, for Theorems~\ref{thm:C} and~\ref{thm:D}, a triple with two entries equal to $1$, from which the required triple is two steps away.
Fact~2 in the proof of Proposition~\ref{prop:classical} was confirmed by a machine enumeration of reduced forms, and Proposition~\ref{prop:classical} itself was checked for $N<1.4\cdot 10^6$.

\subsection*{Scope of the method}
Theorems~\ref{thm:A} and~\ref{thm:B} work because only one binary condition is in play:
the residue of $v$ modulo $3$ in Theorem~\ref{thm:A}, and the equality of the residues of $v$ and $w$ in Theorem~\ref{thm:B}.
In Theorems~\ref{thm:C} and~\ref{thm:D} the additional condition $v\ne 1$ is absorbed by a change of sign.
The sum $p_3+p_5+p_5$ of one triangular and two pentagonal numbers of the same kind, which is also among the $31$ sums of
\cite[Conjecture~1.10]{S15}, asks for a triple in $\cT(24n+5)$ with $v,w\in\{1\}\cup\{5,11,17,\dots\}$;
this is not a consequence of Theorem~\ref{thm:B}, and we have no proof.
The same applies to the tuple $(3,1,3,-1,1,3)$ of \cite{S20}, which needs a triple $(u,v,w)\in\cT(24n+29)$ with $u\ge 3$, $v\equiv 1\pmod 3$ and $w\in\{1\}\cup\{5,11,17,\dots\}$:
Lemma~\ref{lem:large} has no direct analogue for a lower bound on $u$ in $\cT(N)$, and a change of sign does not absorb the condition $u\ne 1$.
By Lemma~\ref{lem:TT} and its proof, the tuple $(3,1,2,6,2,0)$ of \cite{S20} is equivalent to a statement about $\cS(N)$
with the side conditions $|a-b|\ne 2$ and $(a,b)\ne(1,1)$, which we have not been able to treat. In these two cases the alternating sequence of moves reached a triple of the required kind
from every primitive starting triple, for all $n<3000$.

\subsection*{Acknowledgements}
The proofs in this paper were found and written with the assistance of an AI system (Claude, Anthropic),
which also wrote the verification programs. The author is responsible for the content.

\begin{thebibliography}{99}

\bibitem{BI} H. Brandt and O. Intrau, \emph{Tabellen reduzierter positiver tern\"arer quadratischer Formen},
Abh. S\"achs. Akad. Wiss. Math.-Nat. Kl. \textbf{45} (1958), no.~4.

\bibitem{CGQFG} Y. Cao, D. Guo, R. Qiu, R. Feng and X.-S. Gao, \emph{Every nonnegative integer is a sum of a triangular, a pentagonal, and a heptagonal number},
preprint, arXiv:2606.26035, 2026.

\bibitem{Cassels} J. W. S. Cassels, \emph{Rational Quadratic Forms}, Academic Press, London, 1978.

\bibitem{code} Q. Sun, \emph{Verification programs, Lean formalization and data for this paper}, Zenodo, 2026, \url{https://doi.org/10.5281/zenodo.24570762}.

\bibitem{JKS} W. C. Jagy, I. Kaplansky and A. Schiemann, \emph{There are 913 regular ternary forms}, Mathematika \textbf{44} (1997), 332--341.

\bibitem{JP} B. W. Jones and G. Pall, \emph{Regular and semi-regular positive ternary quadratic forms}, Acta Math. \textbf{70} (1939), 165--191.

\bibitem{LK} D. Lorch and M. Kirschmer, \emph{Single-class genera of positive integral lattices}, LMS J. Comput. Math. \textbf{16} (2013), 172--186.

\bibitem{OEIS} OEIS Foundation Inc., \emph{The On-Line Encyclopedia of Integer Sequences}, entry A254574, \url{https://oeis.org/A254574}.

\bibitem{OMeara} O. T. O'Meara, \emph{Introduction to Quadratic Forms}, Springer, Berlin, 1963.

\bibitem{S15} Z.-W. Sun, \emph{On universal sums of polygonal numbers}, Sci. China Math. \textbf{58} (2015), 1367--1396.

\bibitem{S20} Z.-W. Sun, \emph{Universal sums of three quadratic polynomials}, Sci. China Math. \textbf{63} (2020), 501--520.

\end{thebibliography}

\end{document}
